Large language models exhibit two reliability concerns: generating factual errors and succumbing to adversarial perturbations. We hypothesize that calibration quality---the alignment between model confidence and accuracy---serves as a shared internal signal enabling both factuality error detection and adversarial robustness. To test this, we develop a methodology evaluating 12 open-weight LLMs on TruthfulQA (factuality), TextFooler attacks (robustness), and Expected Calibration Error (calibration), then analyzing whether ECE mediates the correlation between factuality and robustness metrics. We present a validated evaluation pipeline and analysis framework. Preliminary results from three models demonstrate pipeline functionality but are insufficient for hypothesis testing; ASR values were simulated and full evaluation remains incomplete. This paper documents our methodology as a foundation for future work connecting these three reliability dimensions.
